\section{工作流程与执行路径}

\subsection{多轮循环框架}

复现任务按图拆分为多轮，每轮独立走 7 步循环：

\begin{verbatim}
第 1 轮：Fig.1（介电球散射截面）  —— 本轮报告
第 2 轮：Fig.2（多极分解）        —— 复用 mie_theory.py
第 3 轮：Grahn 映射验证            —— 复用 code/
第 4 轮：Fig.3 Fano（可选，延后） —— 需求解器或降级
\end{verbatim}

每轮内部的 7 步为：\textbf{01} PDF 预处理（提取参数/公式/图数据）$\to$
\textbf{02} 形式化（写 formalization yaml）$\to$
\textbf{03} 理论笔记（对标讲义 §9--§12）$\to$
\textbf{04} 实现（code + TDD 测试）$\to$
\textbf{05} 运行收集数据 $\to$
\textbf{06} 物理验证（3 层）$\to$
\textbf{07} 分析与报告。

\subsection{四人工作模式}

按 repro-plan-v2.md §2.4 用户拍板的执行方式：

\begin{itemize}
  \item 多模态（读论文图/图表数据点）：codex；
  \item 普通写代码：Claude 内置子 agent（general-purpose/Explore）；
  \item 高能力审查：用户开 claude-opus 独立对话（本报告的表2 blocker 解阻即由
        用户修复 + 主 agent 独立复核完成）；
  \item 物理判断、归因、gate 裁决：Claude 主循环；
  \item 对抗性独立审查：\textbf{W-sub 子 agent}（本报告第 7 节）。
\end{itemize}

\subsection{四个 gate}

\begin{table}[H]
\centering
\caption{四个人工 gate}
\begin{tabular}{clp{7cm}}
\toprule
\textbf{gate} & \textbf{触发点} & \textbf{用户核对内容} \\
\midrule
① 参数 gate & step02 末 & 参数、单位、范围（$\varepsilon_r=6.25$、$2a/\lambda$ 范围） \\
② spec gate & formalization 后 & 物理形式化 spec 与论文物理问题一致 \\
③ 公式 gate & step04/05 末 & Mie 系数 $a_n,b_n$ 对 B\&H 教材核对 \\
④ 误差 gate & step07 末 & \textbf{看量化误差数字，不接受"看起来一致"} \\
\bottomrule
\end{tabular}
\end{table}

\subsection{关键换算：横轴与尺寸参数}

Alaee 横轴是 \textbf{$2a/\lambda$}，Mie 系数用 Bohren--Huffman 尺寸参数 $x=ka$。
host 为空气（$\varepsilon_d=1$）时：
\begin{equation}
  x = ka = \frac{2\pi a}{\lambda} = \pi\cdot\frac{2a}{\lambda}.
  \label{eq:xconv}
\end{equation}
这一 $\pi$ 换算是 gate③ 强调的翻车点：若把 $2a/\lambda$ 直接当 $x$ 用，
Mie 共振峰位将整体偏移 $\pi$ 倍。锚点验证（step04 实测，$n=2.5$）：

\begin{table}[H]
\centering
\caption{共振峰位锚点（实测值与换算校验）}
\begin{tabular}{lccc}
\toprule
\textbf{多极} & \textbf{Mie 系数} & \textbf{$2a/\lambda$ 峰位} & \textbf{$x$ 峰位} \\
\midrule
ED 电偶极 & $a_1$ & 0.500 & $\pi/2\approx1.571$（Fröhlich 共振） \\
MD 磁偶极 & $b_1$ & 0.385 & 1.209（$ka\cdot n\approx\pi$ 球内驻波） \\
EQ 电四极 & $a_2$ & 0.647 & 2.033 \\
MQ 磁四极 & $b_2$ & 0.543 & 1.705 \\
\bottomrule
\end{tabular}
\end{table}
